Criptografia. A palavra pode sugerir uma matéria reservada a matemáticos e
cientistas da computação: obscura, poderosa e, para quem lhe conhece as peças,
elegante. Na verdade, muitos dos seus conceitos fundamentais estão ao alcance
de qualquer leitor disposto a acompanhar o raciocínio passo a passo.

Uma cripto-moeda usa criptografia para definir e transferir a posse de fundos.
Para impedir que a mesma moeda seja gasta duas vezes ou que novas moedas sejam
criadas arbitrariamente, as transacções são registadas numa lista de blocos
(\emph{blockchain}): um livro de contas público e distribuído que terceiros
podem verificar \cite{Nakamoto_bitcoin}. À primeira vista, essa verificabilidade
parece exigir que remetentes, destinatários e montantes fiquem expostos. Monero
mostra que não. É possível ocultar esses dados e, ainda assim, permitir que a
rede confirme a validade das transacções \cite{cryptoNoteWhitePaper}.

Para os leitores experientes: a rede principal usa criptografia sobre a curva
Edwards25519 \cite{Bernstein2008}, endereços de utilização única, assinaturas
em anel CLSAG, compromissos de Pedersen e provas de intervalo Bulletproof+.
Cada anel tem 16 membros e as saídas incluem etiquetas de vista. RandomX
protege a prova de trabalho; Dandelion++ intervém na difusão inicial de
transacções, não nas regras de consenso. Este livro constrói a intuição e a
notação necessárias antes de reunir estas peças.

A segunda edição original descreve o protocolo v12. A edição portuguesa está
a ser revista progressivamente para o protocolo v16; por isso, os capítulos de
protocolo que ainda não tenham sido actualizados devem ser lidos como uma
descrição histórica da v12. O estado técnico acima foi conferido no código
Monero da revisão
\texttt{3646f648db57f60cca86430e25a635d19fa9b92a}, de 14 de Agosto de 2026.
%For our experienced readers: Monero is a standard one-dimensional distributed acyclic graph (DAG) cryptocurrency blockchain \cite{Nakamoto_bitcoin} where transactions are based on elliptic curve cryptography using curve Ed25519 \cite{Bernstein2008}, transaction inputs are signed with Schnorr-style multilayered linkable spontaneous anonymous group signatures (MLSAG) \cite{MRL-0005-ringct}, and output amounts (communicated to recipients via ECDH \cite{Diffie-Hellman}) are concealed with Pedersen commitments \cite{maxwell-ct} and proven in a legitimate range with Bulletproofs \cite{Bulletproofs_paper}. Much of the first part of this report is spent explaining these ideas.
